% [原文 p.167]
\chapter{全纯辛流形}

称复流形 $(M,J)$ 与其上的非退化全纯 $2$ 次微分形式 $\sigma$ 的组 $(M,J,\sigma)$ 为\emph{全纯辛流形}。这是辛流形在复几何中的对应物。紧复流形具有全纯辛形式，是远比紧实流形具有辛形式更严格的条件，已知的具体例子极为有限。另一方面，若去掉紧致性假定，则存在各种全纯辛流形，其中一些是作为场论中的模空间而出现的重要空间。

\section{不可约全纯辛流形}

考虑辛流形在复几何中的类似物，便自然引出下述定义：

\begin{definition}{10.1.1}
  $2n$ 维复流形 $X$ 与其上的全纯 $2$ 次微分形式 $\sigma\in H^0(\Omega_X^2)$（且 $\sigma^n$ 处处不为零）的组 $(X,\sigma)$，称为\emph{全纯辛流形}（holomorphic symplectic manifold）。当全纯辛流形满足
  \begin{enumerate}
    \item[(i)] $X$ 紧致且为 Kähler 流形；
    \item[(ii)] $X$ 单连通；
    \item[(iii)] $H^0(\Omega_X^2)$ 由 $\sigma$ 张成
  \end{enumerate}
  这三个条件时，称为\emph{不可约全纯辛流形}（irreducible holomorphic symplectic manifold）。
\end{definition}

% [Tutor Note: 难点] $X$ 为 $2n$ 维时 $\sigma^n$ 是 $(n,n)$ 型的体积形式，故「$\sigma^n$ 处处非零」等价于 $\sigma$ 非退化。

对全纯辛流形 $(X,\sigma)$，由于 $\sigma^n$ 给出 $\Omega_X^n$ 的平凡化，$X$ 的第 1 Chern 类为零。在第 1 Chern 类为零的流形之中，

% [原文 p.168]
把对任意 $0<p<\dim X$ 都满足 $h^0(\Omega_X^p)=0$ 的流形称为\emph{狭义 Calabi--Yau 流形}。不可约全纯辛流形与环面以及狭义 Calabi--Yau 流形一起，构成第 1 Chern 类为零的 Kähler 流形的基本构件：

% [Tutor Note: 背景] 条件 $h^0(\Omega_X^p)=0$（$0<p<\dim X$）排除了环面与超凯勒因子的混合，使分解中的因子唯一到有限覆盖。

\begin{theorem}{10.1.2}
  （Bogomolov 分解 [24]）第 1 Chern 类为零的紧 Kähler 流形 $X$，经取有限次不分歧覆盖，可替换为复环面 $T$、不可约全纯辛流形 $X_1,\dots,X_r$ 与狭义 Calabi--Yau 流形 $K_1,\dots,K_l$ 的直积 $T\times X_1\times\cdots\times X_r\times K_1\times\cdots\times K_l$。
\end{theorem}

\takeaway{全纯辛形式使 $\Omega_X^n$ 平凡、$c_1=0$；不可约全纯辛流形是 $c_1=0$ 的 Kähler 流形三大构件之一。}

\section{四元数}

1805 年生于爱尔兰 Dublin 的 Hamilton 创始了以他命名的 Hamilton 形式分析力学，在数学与物理学上留下了巨大足迹。然而，他本人认为是生平最大业绩的，却是四元数的发现。据有名的轶事，Hamilton 长年苦心试图在 $\R^3$ 上赋予除环结构，1843 年 10 月 16 日在 Dublin 走过 Broom 桥时，他发觉若考虑 $\R^4$ 而非 $\R^3$，便可定义满足乘法交换性以外的全部除环公理的代数，于是兴奋地把基本关系式
\begin{equation}
  i^2=j^2=k^2=ijk=-1 \tag{10.2.1}
\end{equation}
刻在桥上。为纪念这段故事，这座桥上嵌有刻着下述文字的盾形石板：
\begin{quote}
  Here as he walked by\\
  on the 16th of October 1843\\
  Sir William Rowan Hamilton\\
  in a flash of genius discovered\\
  the fundamental formula for\\
  quaternion multiplication\\[4pt]
  $i^2=j^2=k^2=ijk=-1$\\[4pt]
  \& cut it on a stone of this bridge.
\end{quote}

由 $3$ 个元 $i,j,k$ 生成、具有基本关系式 (10.2.1) 的 $\R$ 上结合代数，因 Hamilton 而记作 $\HH$，其元素称为\emph{四元数}（quaternion）。四元数代数 $\HH$ 是 $\R$ 上的 $4$ 维向量空间
\begin{equation}
  \HH=\R\oplus\R i\oplus\R j\oplus\R k \tag{10.2.2}
\end{equation}
并给出\emph{除环}（skew field）最重要的例子。$\R$ 作为 $\HH$ 的中心自然嵌入 $\HH$，而 $\C$ 到 $\HH$ 的自然嵌入

% [原文 p.169]
\begin{equation}
  \C\to\HH,\qquad a+b\sqrt{-1}\mapsto a+bi \tag{10.2.3}
\end{equation}
的像不是 $\HH$ 的中心。因此，$\C$ 到 $\HH$ 的右作用与左作用不一致，$\HH$ 上具有两种自然的 $\C$ 线性空间结构。用四元数的共轭
\begin{align}
  x=a+bi+cj+dk\mapsto\bar x&:=-\frac{1}{2}(x+ixi+jxj+kxk) \tag{10.2.4}\\
  &=a-bi-cj-dk \tag{10.2.5}
\end{align}
定义四元数的实部、虚部与范数：
\begin{align}
  \Re x&=\frac{1}{2}(x+\bar x)=a, \tag{10.2.6}\\
  \Im x&=\frac{1}{2}(x-\bar x)=bi+cj+dk, \tag{10.2.7}\\
  |x|^2&=x\bar x=a^2+b^2+c^2+d^2. \tag{10.2.8}
\end{align}
% [Tutor Note: 难点] 把 $i^{-1}=-i$ 等代入 $-\frac12(x+ixi+jxj+kxk)$，四个 $i,j,k$ 项两两相消，即得 $a-bi-cj-dk$。

与复数的情形不同，四元数代数上代数方程的解的个数未必与次数一致。实际上 $i,j,k$ 全都满足代数方程
\begin{equation}
  x^2=-1, \tag{10.2.9}
\end{equation}
而 (10.2.9) 的解甚至不是有限个，其解集构成球面
\begin{equation}
  S^2=\{x\in\Im\HH\mid |x|=1\}. \tag{10.2.10}
\end{equation}

对任意 $x,y\in\HH$ 有
\begin{equation}
  \overline{xy}=\bar y\bar x, \tag{10.2.11}
\end{equation}
因此四元数的范数是乘性的；
\begin{equation}
  |xy|=|x||y|. \tag{10.2.12}
\end{equation}
于是 $\HH$ 内的单位球面
\begin{equation}
  S^3=\{x\in\HH\mid |x|=1\} \tag{10.2.13}
\end{equation}
构成群。作为与部分域 $\C\subset\HH$ 的公共部分所得的子群
\begin{equation}
  S^1=S^3\cap\C \tag{10.2.14}
\end{equation}
不是正规子群。由映射
\begin{equation}
  \varphi\colon\HH\to\HH,\qquad x\mapsto xi\bar x \tag{10.2.15}
\end{equation}
给出的 $S^3$ 的像是
\begin{equation}
  S^2=\{x\in\Im\HH\mid |x|=1\}. \tag{10.2.16}
\end{equation}

% [原文 p.170]
并且，对任意 $x,y\in\HH$，$\varphi(x)=\varphi(y)$ 与「存在 $z\in S^1$ 使 $x=yz$」等价。因此商空间 $S^3/S^1$ 与 $S^2$ 同胚，映射 $\varphi\colon S^3\to S^2$ 给出 \emph{Hopf 纤维丛}（Hopf fibration）。

\takeaway{四元数是 $\R^4$ 上的除环；单位四元数群经 $x\mapsto xi\bar x$ 给出 Hopf 纤维丛 $S^3\to S^2$。}

四元数环通过
\begin{equation}
  x=\alpha+j\beta\mapsto\begin{pmatrix}\alpha&-\bar\beta\\ \beta&\bar\alpha\end{pmatrix},\qquad \alpha,\beta\in\C \tag{10.2.17}
\end{equation}
自然嵌入 $2$ 行 $2$ 列的复矩阵环 $M_2(\C)$。该嵌入用 Pauli 矩阵
\begin{equation}
  \sigma_1=\begin{pmatrix}0&1\\1&0\end{pmatrix},\quad
  \sigma_2=\begin{pmatrix}0&-\sqrt{-1}\\\sqrt{-1}&0\end{pmatrix},\quad
  \sigma_3=\begin{pmatrix}1&0\\0&-1\end{pmatrix} \tag{10.2.18}
\end{equation}
可写作
\begin{equation}
  a+bi+cj+dk\mapsto a+b\sqrt{-1}\sigma_3+c\sqrt{-1}\sigma_1+d\sqrt{-1}\sigma_2. \tag{10.2.19}
\end{equation}
同样地，$n$ 行 $n$ 列的四元数矩阵环经嵌入
\begin{equation}
  X+jY\mapsto\begin{pmatrix}X&-\bar Y\\ Y&\bar X\end{pmatrix} \tag{10.2.20}
\end{equation}
与 $2n$ 行 $2n$ 列复矩阵环的部分环等同；
\begin{equation}
  M_n(\HH)=\{g\in M_{2n}(\C)\mid J_{2n}g=\bar gJ_{2n}\}. \tag{10.2.21}
\end{equation}
四元数矩阵环与酉群的公共部分
\begin{equation}
  U(n,\HH):=M_n(\HH)\cap U(2n)\subset M_{2n}(\C) \tag{10.2.22}
\end{equation}
称为\emph{超酉群}（hyperunitary group）。它与复辛群
\begin{equation}
  \Sp(2n,\C)=\{g\in\GL(2n,\C)\mid g^TJ_{2n}g=J_{2n}\} \tag{10.2.23}
\end{equation}
的紧实形式（compact real form）即紧辛群（compact symplectic group）
\begin{equation}
  \Sp(n):=U(2n)\cap\Sp(2n,\C) \tag{10.2.24}
\end{equation}
一致。此外，紧辛群也称为酉辛群（unitary symplectic group），记作 $USp(n)$。另一方面，复 Lie 群 $\Sp(2n,\C)$ 的分裂实形式（split real form）是实辛群
\begin{equation}
  \Sp(2n,\R)=\{g\in\GL(2n,\R)\mid g^TJ_{2n}g=J_{2n}\}. \tag{10.2.25}
\end{equation}
由定义立即得到

% [原文 p.171]
\begin{equation}
  \Sp(1)=S^3 \tag{10.2.26}
\end{equation}
成立。此外，该群作为 $M_2(\C)$ 的子集与 $SU(2)$ 一致。对 $g\in S^3$，定义 $\R$ 线性映射 $\rho(g)\colon\HH\to\HH$ 为
\begin{equation}
  \rho(g)(x)=gx\bar g, \tag{10.2.27}
\end{equation}
则它把 $\Im\HH$ 的元素映入 $\Im\HH$，并保持 $\Im\HH$ 的正定内积
\begin{equation}
  \langle x,y\rangle=\Re(\bar xy). \tag{10.2.28}
\end{equation}
因此，$\rho$ 定义从 $S^3$ 到 $O(3)$ 的同态。该同态 $\rho$ 的核为 $\{\pm 1\}\subset S^3$，像是 $SO(3)$，而 $S^3$ 与 $SO(3)$ 的普遍覆叠群 $Spin(3)$ 等同。

% [Tutor Note: 难点] 该二重覆盖 $S^3\cong SU(2)\to SO(3)$ 即 $Spin(3)$，$\pi_1(SO(3))=\Z/2$ 即由此得到。

\section{超凯勒流形}

把 Kähler 流形的定义从复数扩展到四元数，便得到下述定义：

\begin{definition}{10.3.1}
  $4n$ 维实流形 $M$ 与 $M$ 上 $3$ 个复结构 $I,J,K\in\End(TM)$ 的组 $(M,I,J,K)$，若满足
  \begin{equation}
    I^2=J^2=K^2=IJK=-\id_{TM}, \tag{10.3.1}
  \end{equation}
  则称为\emph{超复流形}（hypercomplex manifold）。超复流形 $(M,I,J,K)$ 与其上的 Riemann 度量 $g$ 的组 $(M,I,J,K,g)$，若 $g$ 对 $3$ 个复结构全都成为 Kähler 度量，则称为\emph{超凯勒流形}（hyperkähler manifold）。
\end{definition}

$4n$ 维 Riemann 流形具有超凯勒流形结构，与它的和乐群含于 $\Sp(n)$ 是等价的。给定超凯勒流形 $(M,I,J,K,g)$ 时，把 $3$ 个复结构 $I,J,K$ 对应的 Kähler 形式分别记作 $\omega_I,\omega_J,\omega_K$。此时，
\begin{equation}
  \sigma_I:=\omega_J+\sqrt{-1}\omega_K \tag{10.3.2}
\end{equation}
给出复流形 $(M,I)$ 上的全纯辛形式。依据 Yau [145] 对 Calabi 猜想的解决，给定紧 Kähler 流形 $X$ 上的全纯辛形式时，在同一上同调类中替换 Kähler 形式，便可使 $X$ 成为超凯勒流形。

\takeaway{超凯勒流形是三个复结构下的 Kähler 流形，其 $\sigma_I=\omega_J+\sqrt{-1}\omega_K$ 给出全纯辛形式。}

对不可约全纯辛流形，各种强断言成立：变形空间为 $b_2(X)-2$ 维，且必为无障碍的（unobstructed）（[25]）。无

% [原文 p.172]
限小 Torelli 问题答案是肯定的，周期映射是满射（[67]）。由这些事实可知，任意超凯勒流形都可经变形成为射影代数簇。此外，整体 Torelli 型定理也成立（[133]）。更进一步，作为显著的事实，有松下与 Hwang 的如下结果：

\begin{theorem}{10.3.2}
  （[68], [88], [89]）设 $X$ 为射影不可约全纯辛流形，$f\colon X\to B$ 为到射影正规簇的固有满射（proper surjection），且 $0<\dim B<\dim X$，$f$ 的一般纤维连通。此时 $f$ 的一般纤维是 Lagrange 环面，底空间 $B$ 成为射影空间。
\end{theorem}

这里所谓 Lagrange 环面，指中间维数的阿贝尔簇，且全纯辛形式在其上的限制为零。这种意义上的 Lagrange 环面纤维丛常称为代数完全可积系。上述定理断言：射影不可约辛流形若具有纤维空间的结构，则必成为代数完全可积系。

这样，由 $X$ 的不可约全纯辛性可导出各种强结论，这正说明不可约全纯辛性是很强的条件。实际上，构造不可约全纯辛流形的例子极为困难，已知的只有 K3 曲面上层的模空间，或与一般化的 Kummer 流形相关的例子。

作为 K3 曲面上稳定层的模空间来构造全纯辛流形 [99]，是「作为曲线上稳定向量丛的模空间来构造实辛流形」这一构造在全纯辛流形中的类似。模空间的切空间由 $\operatorname{Ext}^1(\mathcal{E},\mathcal{E})$ 给出，其上的配对不是杯积与 Poincaré 对偶性，而是米田积与 Serre 对偶性合成的
\begin{equation}
  \operatorname{Ext}^1(\mathcal{E},\mathcal{E})\times\operatorname{Ext}^1(\mathcal{E},\mathcal{E})
  \to\operatorname{Ext}^2(\mathcal{E},\mathcal{E})\cong\operatorname{Ext}^0(\mathcal{E},\mathcal{E})^*\cong\C \tag{10.3.3}
\end{equation}
稳定层模空间的典型例子，是 K3 曲面上点的 Hilbert 概型。

取阿贝尔曲面 $A$，考虑 $A$ 的加法所定的态射 $A^n\to A$。由加法的可换性，它诱导从对称积出发的态射 $S^nA\to A$。把它与 Hilbert--Chow 态射 $A^{[n]}\to S^nA$ 合成，便得到从 $A$ 的 $n$ 点 Hilbert 概型 $A^{[n]}$ 出发的态射 $A^{[n]}\to A$。该态射下 $0\in A$ 的原像 $K^{(n-1)}A$ 称为 $(n-1)$ 阶一般化 Kummer 流形（$(n-1)$-st generalized Kummer variety），它是 $2n-2$ 维的不可约全纯辛流形（[19]）。

\section{超凯勒商}

与紧致例子之稀少形成对照，非紧的全纯辛流形与超凯勒流形的例子却大量存在。事实上，对任意复流形，其全纯余切丛都给出非紧全纯辛流形的例子。构造超凯勒流形的例子比构造全纯辛流形的例子更为困难，但通过下述辛商在超凯勒流形中的类比，可以构造出许多例子。

% [原文 p.173]
设 Lie 群 $G$ 作用于超凯勒流形 $(M,g,I,J,K)$。若该作用对 $3$ 个 Kähler 形式 $\omega_I,\omega_J,\omega_K$ 全都成为 Hamilton 作用，则称为\emph{超 Hamilton 作用}（hyperhamiltonian action）。此时把 $3$ 个动量映射 $\mu_I,\mu_J,\mu_K$ 合并写作
\begin{equation}
  \mu=(\mu_I,\mu_J,\mu_K)\colon M\to\mathfrak{g}^*\otimes\Im\HH \tag{10.4.1}
\end{equation}
并称为\emph{超凯勒动量映射}（hyperkähler moment map）。$G$ 通过对 $\mathfrak{g}^*$ 的余伴随作用与对 $\Im\HH$ 的平凡作用的张量积表示作用于 $\mathfrak{g}^*\otimes\Im\HH$，把该作用的固定点集合记作 $(\mathfrak{g}^*\otimes\Im\HH)^G$。

\begin{theorem}{10.4.1}
  （[63]）设 Lie 群 $G$ 在超凯勒流形 $(M,g,I,J,K)$ 上具有超 Hamilton 作用，其超凯勒动量映射由
  \begin{equation}
    \mu=(\mu_I,\mu_J,\mu_K)\colon M\to\mathfrak{g}^*\otimes\Im\HH \tag{10.4.2}
  \end{equation}
  给出。此时，对任意 $\zeta\in(\mathfrak{g}^*\otimes\Im\HH)^G$，若 $G$ 在 $\mu^{-1}(\zeta)$ 上的作用自由，则
  \begin{equation}
    M/\!/\!/_\zeta G:=\mu^{-1}(\zeta)/G \tag{10.4.3}
  \end{equation}
  是 $(\dim M-4\dim G)$ 维的超凯勒流形。
\end{theorem}

% [Tutor Note: 难点] 超凯勒商是辛商的三重叠加：每个动量映射各降 $2\dim G$ 维，三个合起来恰好降 $4\dim G$ 维。

\takeaway{超 Hamilton 作用与超凯勒动量映射给出超凯勒商，其维数为 $\dim M-4\dim G$。}

\section{全纯辛向量空间}

最基本的辛流形是辛向量空间
\begin{equation}
  \Bigl(\C^n,\ \frac{\sqrt{-1}}{2}\sum_{i=1}^n dz_i\wedge d\bar z_i\Bigr), \tag{10.5.1}
\end{equation}
而它在全纯辛几何中的类似物便是\emph{全纯辛向量空间}（holomorphic symplectic vector space）。把 $\HH^n$ 的元素用 $z,w\in\C^n$ 写作 $z+jw$ 时，$3$ 个 Kähler 形式为
\begin{equation}
  \omega_I=\frac{\sqrt{-1}}{2}\sum_{i=1}^n\bigl(dz_i\wedge d\bar z_i+dw_i\wedge d\bar w_i\bigr), \tag{10.5.2}
\end{equation}
\begin{equation}
  \sigma_I=\omega_J+\sqrt{-1}\omega_K=\sum_{i=1}^n dz_i\wedge dw_i. \tag{10.5.3}
\end{equation}

% [Tutor Note: 难点] 写法 $z+jw$ 中的 $j$ 是不作右乘的虚单位，$z,w$ 取复坐标；故 $dz_i\wedge dw_i$ 为 $(2,0)$ 型，$\sigma_I$ 全纯。

\takeaway{$\HH^n\cong\C^{2n}$ 上的 $\sigma_I=\sum dz_i\wedge dw_i$ 给出全纯辛形式，是辛向量空间的全纯类比。}

\section{超环面流形}

阿贝尔群的同态 $\varphi\colon\Z^d\to\Z^n$ 定出 $\R^d/\Z^d$ 的 $(d-n)$ 维部分环面 $K=\ker\varphi\otimes(\R/\Z)$。

% [原文 p.174]
另一方面，$\R^d/\Z^d$ 在 $\C^d$ 上有自然的 Hamilton 作用
\begin{equation*}
  (\alpha_1,\dots,\alpha_d)\colon z=(z_1,\dots,z_d)
  \mapsto t\cdot z=(e^{2\pi\sqrt{-1}\alpha_1}z_1,\dots,e^{2\pi\sqrt{-1}\alpha_d}z_d),
\end{equation*}
其动量映射由
\begin{equation}
  z\mapsto\frac{1}{2}\sum_{i=1}^d|z_i|^2e_i \tag{10.6.1}
\end{equation}
给出。把该 Hamilton 作用限制到 $K$ 上，$K$ 便自然作用于 $\C^d$。关于这个 Hamilton 作用，$\C^d$ 在 $\lambda\in\mathfrak{t}^*$ 处的辛商 $\C^d/\!/_\lambda K$ 便称为辛环面流形。

实环面 $\R^d/\Z^d$ 在辛向量空间 $\C^d$ 上的 Hamilton 作用 $z\mapsto t\cdot z$ 引起 $\R^d/\Z^d$ 在全纯辛向量空间 $T^*\C^d\cong\C^{2d}$ 上的超 Hamilton 作用
\begin{equation}
  t\cdot(z,w)=(t\cdot z,t^{-1}\cdot w), \tag{10.6.2}
\end{equation}
其超动量映射由
\begin{equation}
  \mu_\R(z,w)=\mu_1(z,w)=\frac{1}{2}\sum_{i=1}^d\bigl(|z_i|^2-|w_i|^2\bigr)e_i, \tag{10.6.3}
\end{equation}
\begin{equation}
  \mu_\C(z,w)=(\mu_2+\sqrt{-1}\mu_3)(z,w)=\sum_{i=1}^d z_iw_ie_i \tag{10.6.4}
\end{equation}
给出。对 $\zeta\in\mathfrak{t}^*\otimes\Im\HH$，把上述超 Hamilton 作用限制到 $K$ 上所得的超凯勒商
\begin{equation}
  Y(\varphi,\zeta):=\C^{2d}/\!/\!/_\zeta K \tag{10.6.5}
\end{equation}
称为\emph{超环面流形}（hypertoric manifold，或 toric hyperkähler manifold）。它由 Bielawski--Dancer [22] 引入，并由今野 [78] 与 Hausel--Sturmfels [60] 进一步研究。关于超环面流形的解说，例如可参看 [116]。

% [Tutor Note: 背景] 同一环面作用在 $\C^d$ 上取辛商得辛环面流形，在 $T^*\C^d$ 上取超凯勒商得超环面流形，后者是前者的超凯勒类比。

\takeaway{对同一环面作用，辛商给出辛环面流形，超凯勒商给出超环面流形 $Y(\varphi,\zeta)$。}

\section{箭图簇}

箙（quiver）即有向图，形式地定义为
\begin{itemize}
  \item 顶点（vertex）的集合 $Q_0$；
  \item 箭头（arrow）的集合 $Q_1$；
  \item 把箭头的起点（source）与终点（target）对应起来的映射 $s,t\colon Q_1\to Q_0$
\end{itemize}
的组 $Q=(Q_0,Q_1,s,t)$。

% [Tutor Note: 背景] 原书节名为「箙多様体」；「箙」即 quiver（箭图），故本节标题按其数学内容译作「箭图簇」。

% [原文 p.175]
当箙 $Q$ 的各顶点 $v\in Q_0$ 配上复向量空间 $V_v$ 与 $W_v$ 时，置
\begin{equation}
  V=\bigoplus_{a\in Q_1}\Hom(V_{s(a)},V_{t(a)}),\qquad
  W=\bigoplus_{v\in Q_0}\Hom(W_v,V_v), \tag{10.7.1}
\end{equation}
则 $V\oplus W$ 上 Lie 群
\begin{equation}
  U(V):=\prod_{v\in Q_0}U(V_v) \tag{10.7.2}
\end{equation}
自然地通过
\begin{equation}
  B_a\in\Hom(V_{s(a)},V_{t(a)})\mapsto g_{t(a)}\circ B_a\circ g_{s(a)}^{-1}, \tag{10.7.3}
\end{equation}
\begin{equation}
  i_v\in\Hom(W_v,V_v)\mapsto g_v\circ i_v \tag{10.7.4}
\end{equation}
作用。该作用诱导出全纯辛向量空间 $T^*(V\oplus W)$ 上的超 Hamilton 作用，其超凯勒动量映射为
\begin{equation}
  \mu_\R=\frac{\sqrt{-1}}{2}\Bigl(\sum_{a\in Q_1}\bigl(B_aB_a^*-B_a^*B_a\bigr)
  +i_vi_v^*-j_v^*j_v\Bigr)_{v\in Q_0}, \tag{10.7.5}
\end{equation}
\begin{equation}
  \mu_\C=\Bigl(\sum_{t(a)=v}\varepsilon B_aB_{\bar a}+i_vj_v\Bigr)_{v\in Q_0}. \tag{10.7.6}
\end{equation}
对应的超凯勒商称为\emph{箭图簇}（quiver variety）。它由中岛在研究 ALE 空间上的瞬子时引入，并应用于几何表示论 [105], [106]。

% [Tutor Note: 难点] (10.7.6) 中的 $\varepsilon$ 取自 $\{\pm1\}$，其符号随复结构 $I,J,K$ 的选取而变，原书未在此处说明。

取箙
\begin{equation}
  Q_0=\{0,1,\dots,n\},\qquad Q_1=\{1,2,\dots,n\}, \tag{10.7.7}
\end{equation}
\begin{equation}
  s(i)=i,\quad t(i)=0,\qquad i=1,\dots,n, \tag{10.7.8}
\end{equation}
并取附属于各顶点的向量空间
\begin{equation}
  V_0=\C^2,\quad V_i=\C,\qquad i=1,\dots,n, \tag{10.7.9}
\end{equation}
\begin{equation}
  W_i=0,\qquad i=0,\dots,n \tag{10.7.10}
\end{equation}
时，所得的箭图簇称为\emph{超多角形空间}（hyperpolygon space）。它作为多角形空间在超凯勒几何中的类比由今野 [79] 引入，并已知与 $\PP^1$ 上秩 $2$ 稳定抛物 Higgs 束的模空间同型（[51]）。

% [Tutor Note: 背景] 箭图簇把几何表示论中的 Nakajima 对应与超凯勒商联系起来：取 $W_v=0$ 时所得即 ALE 空间上的瞬子模空间。

\takeaway{箭图配上向量空间后，$U(V)$ 的超凯勒商即箭图簇；特定取材给出超多角形空间。}

\chapterreview{本章把辛几何的复化推广到全纯辛流形，并给出其最重要的例子类。全纯辛形式使 $\Omega_X^n$ 平凡、$c_1=0$，而不可约全纯辛流形与复环面、狭义 Calabi--Yau 流形一起构成 $c_1=0$ 紧 Kähler 流形的基本构件（Bogomolov 分解）。四元数代数 $\HH$ 提供三个复结构的模型，超凯勒流形正是对 $\omega_I,\omega_J,\omega_K$ 全为 Kähler 的流形，$\sigma_I=\omega_J+\sqrt{-1}\omega_K$ 给出全纯辛形式。超凯勒商是辛商在超凯勒情形下的类比，$\C^d$ 与 $T^*\C^d$ 上的同一环面作用分别给出辛环面流形与超环面流形。箭图簇作为 $U(V)$ 的超凯勒商，把几何表示论与超凯勒几何联结起来。}